\section{Sprint 0 : Socle Technique}
\textbf{Objectif} : Mettre en place l'environnement et le pipeline d'ingestion.
\textbf{Réalisations} :
\begin{itemize}
    \item Création de la structure du projet (MVC-like).
    \item Implémentation du wrapper \texttt{ffmpeg\_wrapper.py} pour l'extraction audio/vidéo.
    \item Définition du schéma de données JSON (Dataclasses).
\end{itemize}

\section{Sprint 1 : Transcription (ASR)}
\textbf{Objectif} : Convertir la parole en texte horodaté.
\textbf{Réalisations} :
\begin{itemize}
    \item Intégration de \texttt{faster-whisper}.
    \item Utilisation du modèle "medium" avec quantification int8 pour la performance CPU.
    \item Extraction précise des timestamps (début/fin de phrase).
\end{itemize}

\section{Sprint 2 : Métriques Audio}
\textbf{Objectif} : Analyser "comment" la personne parle.
\textbf{Réalisations} :
\begin{itemize}
    \item \textbf{WPM (Words Per Minute)} : Calculé sur la durée active de parole.
    \item \textbf{Pauses} : Détection des silences > 0.5s via l'énergie RMS (Librosa).
    \item \textbf{Fillers} : Comptage des tics de langage (euh, um) via Regex sur le transcript.
\end{itemize}

\section{Sprint 3 : Analyse Visuelle (Base)}
\textbf{Objectif} : Vérifier la présence et la qualité vidéo.
\textbf{Réalisations} :
\begin{itemize}
    \item Intégration de \texttt{MediaPipe Face Mesh}.
    \item Calcul du "Face Presence Ratio" (pourcentage de frames avec un visage détecté).
    \item Analyse basique de la qualité (Luminosité moyenne, Flou via Laplacien).
\end{itemize}

\section{Sprint 4 : Analyse du Regard (Eye Contact)}
\textbf{Objectif} : Estimer si l'orateur regarde la caméra.
\textbf{Réalisations} :
\begin{itemize}
    \item \textbf{Head Pose Estimation} : Calcul des angles d'Euler (Yaw, Pitch, Roll) via l'algorithme PnP (\texttt{cv2.solvePnP}) sur 6 landmarks faciaux clés (Nez, Menton, Yeux, Bouche).
    \item \textbf{Heuristique Eye Contact} : On considère qu'il y a contact visuel si $|Yaw| < 15^{\circ}$ et $-25^{\circ} < Pitch < 15^{\circ}$ (tolérance augmentée vers le bas pour compenser l'angle d'écran de PC portable).
    \item \textbf{Intégration} : Ajout de la métrique "Eye Contact Ratio" dans le rapport final.
\end{itemize}

\section{Sprint 5 : Analyse Gestuelle (Mains et Posture)}
\textbf{Objectif} : Quantifier l'usage des mains comme support au discours.
\textbf{Réalisations} :
\begin{itemize}
    \item \textbf{Hand Detection} : Intégration de \texttt{MediaPipe Hands}.
    \item \textbf{Hand Visibility Ratio} : Pourcentage de frames où au moins une main est détectée. Permet de savoir si l'orateur est "corps figé" ou expressif.
    \item \textbf{Gesture Intensity} : Calcul de la vitesse moyenne des poignets (normalisée sur une échelle 0-10). Cela permet de détecter l'hyperactivité (distrayante) ou l'absence de mouvement.
\end{itemize}

\section{Sprint "Consolidation" : Visualisation & Exports}
\textbf{Objectif} : Générer les livrables manquants de la roadmap (fichiers bruts, graphiques) pour la démo.
\textbf{Réalisations} :
\begin{itemize}
    \item \textbf{Export Transcript} : Génération automatique de \texttt{transcript.txt} séparé.
    \item \textbf{Visualisation Matplotlib} : Génération d'un graphique "Timeline" (\texttt{timeline.png}) montrant les segments de parole actifs pour visualiser le rythme.
    \item \textbf{Intégration Rapport} : Inclusion automatique des graphiques dans le Markdown.
\end{itemize}
